\begin{frame}{채점 기준 (rubric)}
  \small
  \begin{tabular}{@{}llp{0.44\textwidth}@{}}
    \toprule
    항목 & 비중 & 좋은 답 \\
    \midrule
    문제 정의와 환경 & 15\% & MDP 5요소·차원·지표·베이스라인·리턴 이론
    범위가 명확, 환경 선택이 ``개념 X가 돋보이는 무대'' 논리로 근거됨 \\
    방법론 적절성 & 20\% & Ch9$\sim$15 알고리즘이 환경 특성(연속/이산,
    보상 희소성)에 맞았고, 하이퍼파라미터 근거가 배운 개념으로 설명됨 \\
    \textbf{실험 프로토콜의 정직성} & \textbf{25\%} & 시드 $\geq 3$,
    동일 평가 프로토콜(결정적, 동일 에피소드 수), 베이스라인, 예산과 선택
    절차가 보고서에서 추적 가능 \\
    수식적 정당화 & 20\% & 배운 개념 $\geq 1$개로 ``왜''를 설명했고,
    그 설명이 \textbf{실제로 관측한 숫자와 일치} \\
    \textbf{결과의 정직성과 반성} & \textbf{20\%} & 불안정했던 시드·구간,
    베이스라인 이하의 평가치를 숨기지 않고 Ch9$\sim$15 개념으로 진단 \\
    \bottomrule
  \end{tabular}
  \vskip0.8em
  \begin{block}{정직성 항목이 \textbf{합계 45\%}인 이유}
    이 프로젝트의 목적은 ``좋은 정책을 만든 것''이 아니라 \textbf{그 성능
    숫자를 정직하게 얻는 절차를 실제로 수행한 것}을 증명. 결정적 평가가
    베이스라인을 못 넘겨도 프로토콜이 정직하면 높은 점수, 반대로 숫자가
    예뻐도 \textbf{시드 1개였다면} ``프로토콜의 정직성''에서 크게 감점.
  \end{block}
\end{frame}
